\section{Other grids and completed grids}\label{sec:catalogue}

Table~\ref{tab:catalogue} collects the named uncolorable grids and completed grids, including those that lie outside
the families of Sections~\ref{sec:real}--\ref{sec:roots}, with the standard of proof for each
minimum.  ``Certified'' means the procedure of Section~\ref{sec:methods} terminated;
``$\le$'' means an uncolorable subset of that size is in hand and nothing smaller was excluded;
$^{\mathrm L}$ marks a minimum that is also proved in Lean.  A grid is contained in its
completed grid, so a minimum set of the grid bounds the completed grid from above
(Lemma~\ref{lem:restrict}); it need not be a minimum set of it.

\begin{table}[t]
\centering
\caption{Uncolorable grids and completed grids.  ``grid'': all rays whose coordinates are
coordinate values; ``compl.'': the completed grid, the closure under
third rays of orthogonal pairs.  Notation: $\omega=e^{2\pi i/3}$, $\alpha=(1+\sqrt{-7})/2$,
$\varphi=(1+\sqrt5)/2$, $\psi^3=\psi^2+1$ ($\psi\approx1.4656$), $c=\sqrt[3]{2}$.  Bases: number
of triads inside the minimum (or best) set found.  $^{\mathrm L}$: minimum also proved in Lean
(Section~\ref{sec:methods}).}
\label{tab:catalogue}
\small
\begin{tabular}{llrrrll}
\toprule
Coordinate values & Set & Rays & Pairs & Triads & Minimum KS subset & Bases \\
\midrule
$\{0,\pm1,\pm2\}$ & grid & 49 & 138 & 26 & 31, certified$^{\mathrm L}$ & 17 \\
 & compl. & 109 & 258 & 86 & 31, certified & 17 \\
$\{0,\pm1,\pm\sqrt2\}$ & grid & 49 & 120 & 16 & 33, certified$^{\mathrm L}$ & 16 \\
 & compl. & 145 & 336 & 112 & 33, certified & 20 \\
$\{0,\pm1,\pm\sqrt{-2}\}$ & grid & 49 & 120 & 16 & 33, certified$^{\mathrm L}$ & 16 \\
 & compl. & 145 & 336 & 112 & 33, certified & 20 \\
$\{0,\pm1,\pm\omega,\pm\omega^2\}$ & grid & 57 & 174 & 22 & 33, certified$^{\mathrm L}$ & 14 \\
 & compl. & 165 & 390 & 130 & 33, certified & 14 \\
$\{0,\pm1,\pm i,\pm(1+i)\}$ & grid & 127 & 465 & 51 & 33, certified & 16 \\
 & compl. & 583 & 1521 & 507 & $\le33$ & \\
$\{0,\pm1,\pm\alpha,\pm\alpha^{*}\}$ & grid & 145 & 522 & 42 & 43, certified$^{\mathrm L}$ & 23 \\
 & compl. & 685 & 1746 & 582 & \heegnercompletedmin & \\
$\{0,\pm1,\pm\psi,\pm\psi^2\}$ & grid & 109 & 348 & 40 & 47, certified$^{\mathrm L}$ & 26 \\
 & compl. & \superfourcompleted & \superfourmin \\
$\{0,\pm1,\pm\psi\}$ & compl. & \supercompleted & \supermin \\
$\{0,\pm1,\pm\varphi\}$ & compl. & 205 & 498 & 166 & \goldenmin & \goldenbases \\
$\{0,\pm1,\pm c,\pm c^2\}$ & compl. & \cubiccompleted & \cubicmin \\
\bottomrule
\end{tabular}
\end{table}

\paragraph{Completion.}  The last three rows come from a different construction.  The grids
of $\{0,\pm1,\pm\varphi\}$ (49 rays, 10 triads), $\{0,\pm1,\pm\psi\}$ (49 rays, 10 triads) and
$\{0,\pm1,\pm c,\pm c^2\}$ (109 rays, 28 triads) are colorable; uncolorability appears only
after completion, which adds rays that no coordinate triple names.  The completed grids of
$\{0,\pm1,\pm\sqrt3\}$ and $\{0,\pm1,\pm\sqrt5\}$ (133 rays, 282 pairs, 94 triads each) remain
colorable.  Since $\sqrt5$ and $\varphi$ generate the same field, this is the plainest evidence
that the list of coordinate values, and not the field, is the operative choice.  We have no account of why
completion suffices in one case and not in the other.  (Earlier versions said that completion
``generates the reciprocal $1/\varphi$''; that explains nothing, since
$[\varphi:1:0]=[1:1/\varphi:0]$ is already in the grid.)  For these completed grids and for the other
completions marked ``$\le$'' the certifier did not terminate, and the entries are upper bounds:
a 51-ray subset for the golden ratio (118 pairs, 28 bases; coordinates in
\texttt{golden\_51\_rays.txt}; version~9 claimed a certified minimum of 52, which is wrong); a
60-ray subset for the cube root of two (144 pairs, 27 bases); and 47 for both supergolden
completions.

\paragraph{The supergolden grid.}  Let $\psi$ be the real root of $\psi^3=\psi^2+1$.  The
grid of $\{0,\pm1,\pm\psi,\pm\psi^2\}$ is uncolorable without any completion, with certified
minimum 47 (110 pairs, 26 bases); it is one of the eight points of Table~\ref{tab:threesymbol}.
A 47-ray set over this field was first communicated to us by J.~L.~Gonzalez, extracted from the
completed grid of $\{0,\pm1,\pm\psi\}$.  It refuted the classification of versions 1--8, because
no product of two coordinate values of $\{0,\pm1,\pm\psi\}$ has squared modulus~2 and the field contains no roots of unity other
than $\pm1$: the relation at work is the minimal polynomial of $\psi$ itself.  The conclusion we
now draw is that any relation among the coordinate values that makes enough inner products vanish may do.
Theorem~\ref{thm:threesymbol} lists all such relations for three real coordinate values; for more
coordinate values we have no finite classification.

\begin{table}[t]
\centering
\caption{The minimum sets of the certified grids and the 33-ray set of Sch\"utte (Conway
and Kochen), which is critical but not a minimum set of its grid.  Degree sequences are written
$d^{m}$ for $m$ rays orthogonal to $d$ others.  ``Indisp.'': number of bases whose ``exactly
one'' rule cannot be dropped (all orthogonal-pair exclusions kept) without the set becoming
colorable.  Every set is ray-critical.}
\label{tab:profiles}
\small
\begin{tabular}{lrrrll}
\toprule
Set & Rays & Pairs & Bases & Degree sequence & Indisp. \\
\midrule
$\CK$ & 31 & 71 & 17 & $3^{4}4^{14}5^{8}6^{3}8^{2}$ & 17 of 17 \\
Sch\"utte / Conway--Kochen 33 & 33 & 76 & 20 & $3^{8}4^{13}6^{10}8^{2}$ & 20 of 20 \\
Peres 33 & 33 & 72 & 16 & $4^{30}8^{3}$ & 13 of 16 \\
$\sqrt{-2}$, 33 & 33 & 72 & 16 & $4^{30}8^{3}$ & 13 of 16 \\
Eisenstein 33 & 33 & 78 & 14 & $4^{18}5^{12}8^{3}$ & 13 of 14 \\
$\Q(\sqrt{-7})$, 43 & 43 & 107 & 23 & $4^{16}5^{20}6^{4}8^{2}10^{1}$ & 18 of 23 \\
Supergolden 47 & 47 & 110 & 26 & $3^{14}4^{12}5^{8}6^{6}7^{4}8^{1}9^{2}$ & 26 of 26 \\
\bottomrule
\end{tabular}
\end{table}

Table~\ref{tab:profiles} gives the basic invariants of the minimum sets.  The three 33-ray sets
of Conway--Kochen, Peres and the Eisenstein grid have pairwise non-isomorphic orthogonality
graphs (VF2, and visibly from the degree sequences).  The Eisenstein set is the 33-ray set of
Cabello~\cite{Cabello2025simplest}, or one of its three images in the grid
(Computed fact~\ref{fact:eisenstein}); we have matched the invariants he reports (33 rays, 14
bases), and the explicit bijection has not been recorded in the repository.

Whether the 43-ray and 47-ray sets, or the 45-ray family of Theorem~\ref{thm:circle}, appear in
the published catalogues we do not know.  A proper comparison against the large automated
catalogues of three-dimensional KS sets has not been done, and we make no claim that these
configurations are new.
